\chapter{API sebagai Produk dan Ekonomi API}

\section*{Tujuan Pembelajaran}
\begin{itemize}
  \item Menjelaskan konsep API sebagai produk dan perbedaannya dengan API sebagai proyek integrasi.
  \item Mengklasifikasikan API menurut audiensnya: privat, mitra, dan publik.
  \item Menyusun inventaris integrasi dan API yang berjalan di organisasi.
\end{itemize}

\section{Pendahuluan}

% TODO
% Cakupan: biaya integrasi Medikara tinggi karena tiap integrasi dibangun sebagai proyek sekali
% jadi oleh vendor berbeda, dan tidak ada yang tahu pasti berapa banyak integrasi yang berjalan.

\section{Ekonomi API}

% TODO
% Cakupan: arti ekonomi API; temuan Benzell, Hersh, dan Van Alstyne bahwa API membalik batas
% perusahaan dan berkaitan dengan pertumbuhan; temuan Postman 2025 State of the API (adopsi
% API-first, pendapatan dari API); posisi Indonesia melalui SNAP dan SATUSEHAT secara singkat.

\section{API sebagai Produk}

% TODO
% Cakupan: pengguna API, nilai yang ditawarkan, pemilik produk, dan peta jalan; perbedaan dengan
% API sebagai proyek; pendekatan API-first; API sebagai produk internal maupun eksternal.
% Kasus Medikara: API jadwal dokter yang dibuat ulang oleh tiap vendor.

\section{Audiens API: Privat, Mitra, dan Publik}

% TODO
% Cakupan: klasifikasi API privat, mitra, dan publik; konsekuensi pada keamanan, dokumentasi,
% kontrak, dan dukungan; perpindahan API dari privat ke mitra.
% Kasus Medikara: API pendaftaran pasien, API hasil laboratorium, dan API katalog layanan.

\section{Integrasi Point-to-Point dan Utang Integrasi}

% TODO
% Cakupan: pertumbuhan jumlah sambungan point-to-point seiring bertambahnya sistem; biaya
% tersembunyi berupa pemeliharaan, ketergantungan vendor, dan kegagalan berantai; konsep utang
% integrasi.
% Kasus Medikara: integrasi tiap rumah sakit ke BPJS Kesehatan dan SATUSEHAT yang dibangun terpisah.

\section{Model Bisnis Berbasis API}

% TODO
% Cakupan: contoh perusahaan yang produk utamanya API di sektor pembayaran dan komunikasi; contoh
% layanan pembayaran berbasis SNAP di Indonesia; contoh di sektor kesehatan seperti layanan
% laboratorium; monetisasi dibahas rinci di Bab 12.

\section{Menyusun Inventaris Integrasi dan API}

% TODO
% Cakupan: kolom inventaris (sistem asal, sistem tujuan, data, pemilik, cara pertukaran, audiens,
% frekuensi, masalah); sumber informasi berupa wawancara, dokumentasi, konfigurasi, dan kontrak
% vendor; cara memeriksa kelengkapan inventaris.
% Kasus Medikara: contoh baris inventaris untuk SIMRS, ERP, Medikara+, BPJS Kesehatan, dan
% SATUSEHAT.
% Pendalaman teknis: menelusuri katalog API pada SNAP Developer Site milik ASPI dan membandingkan
% cara katalog mendeskripsikan produk API.

\section{Aktivitas: Inventaris Integrasi dan API Saat Ini}

% TODO
% Tujuan aktivitas: memiliki daftar lengkap integrasi dan API yang berjalan sebagai dasar C2.
% Peran dalam tim: anggota berlatar non-teknis mewawancarai pemilik proses tentang kegunaan dan
% masalah tiap integrasi; anggota berlatar teknis menelusuri dokumentasi, konfigurasi, atau kontrak
% vendor.
% Langkah: mulai dari daftar pertukaran data Bab 1; isi templat inventaris; klasifikasikan audiens;
% tandai integrasi paling mahal atau paling rawan; validasi bersama pemilik proses.
% Luaran dan tonggak capstone: inventaris integrasi dan API, masuk ke C2.
% Berkas sources yang direncanakan: sources/chapter04/inventaris-integrasi.csv,
% sources/chapter04/panduan-wawancara-pemilik-proses.md,
% sources/chapter04/contoh-inventaris-medikara.csv

\section{Ringkasan}

% TODO

\section{Latihan}

\begin{enumerate}
  \item Jelaskan dengan contoh mengapa integrasi \textit{point-to-point} yang jumlahnya terus bertambah menaikkan biaya organisasi.
  \item Tentukan apakah API pendaftaran pasien Medikara sebaiknya privat, mitra, atau publik, lalu tuliskan alasannya.
  \item Sebutkan tiga ciri yang membedakan API yang dikelola sebagai produk dari API yang dibangun sebagai proyek sekali jadi.
\end{enumerate}
